\section{Introduction}

Strong gravitational lenses constrain galaxy mass profiles, dark-matter
structure, and cosmological distances \citep[e.g.][]{Treu2010,Collett2015}.
Building large samples requires scalable discovery methods, including
convolutional neural networks (CNNs). Population inference also requires a
selection function: a model of whether a lens is physically detectable and
whether a specified search procedure finds it
\citep{Sonnenfeld2023,Sonnenfeld2024}. Injection--recovery measures the
response after synthetic lensed light is added to survey images. Without a
separate detectability condition, its recovery fraction combines physical
visibility with finder response.

CNN lens finders perform well on simulated benchmarks and labelled
observational samples \citep[e.g.][]{Petrillo2017,Lanusse2018,Jacobs2019,
Canameras2021,Storfer2024}, but their calibration depends on the simulations
used. Some pipelines inject parametric S\'{e}rsic sources lensed by idealised
mass models \citep[e.g.][]{Herle2024}; others use images of observed galaxies
as source-plane templates \citep[e.g.][]{Canameras2021,Rojas2022,Schuldt2023}.
Cross-pipeline results combine differences in target selection, lens models,
source distributions, image processing, and source structure
\citep{Jaelani2024,Canameras2024}. A classifier two-sample test
asks whether a model can tell which of two samples an object came from using
only a chosen representation \citep{LopezPazOquab2017}. In strong-lens work,
self-supervised survey representations have supported lens discovery
\citep{Stein2022}, while domain-adaptation methods have explicitly addressed
differences between simulated and observed lens data \citep{Parul2025}.
A linear probe is an \(L_2\)-regularised logistic regression fitted to frozen
features; its receiver-operating-characteristic area under the curve (AUC)
measures how well the two samples are separated on held-out data.

Injection recovery measures how often a finder responds to a specified
simulated population. A complementary representation-level comparison asks
how clearly observed and injected populations separate in a chosen feature
space. We therefore use a linear probe of the operational detector's pooled
penultimate features as a detector-specific diagnostic. It reuses the trained
representation to evaluate a held-out observed--injected readout for the fixed
finder. We
apply the same dependency-aware procedure to self-supervised survey features,
generic ImageNet features, and elementary pixel statistics to map how broadly
separation appears across representations. Comparing the four protocols shows
whether separation is specific to the lens-trained network or also appears in
unrelated learned and pixel-level representations.

We interpret the comparison using the existing lens-classification score,
native-pixel summaries, and matched observed--bare--injected controls.
We also measure how the fixed detector responds to the injection grid,
source-photon sampling, source construction, and source brightness. Together,
these analyses distinguish
population-level representation differences from responses to controlled
construction choices.
